\section{Introduction}
\label{sec:intro}

Security operations centres (SOCs) triage large alert volumes with limited analyst
capacity~\cite{vielberth2020soc,tariq2025alertfatigue}, and field studies document
recurring process problems: false-positive-dominated alarms, mismatched expectations
between analysts and management, and burnout~\cite{alahmadi2022falsepositives,kokulu2019soc,sundaramurthy2015burnout}.
When an external body supervises many SOCs---for example a national authority for
critical infrastructure---it must judge periodically whether each monitored
organization actually investigated, escalated and closed its work properly. The
instruments available are mostly maturity self-assessments such as
SOC-CMM~\cite{vanos2016soccmm} and CTI-SOC2M2~\cite{schlette2021ctisoc2m2}, metric
catalogues~\cite{forsberg2023metrics,agyepong2023analyst}, and manual examination of
records. Self-assessment reports capability, not conduct; manual examination does not
scale across many organizations and periods.

Record-based assessment of incident handling is not new. Palma et al.\ assess the
compliance of an incident-management process with a reference model derived from
ITIL and ISO/IEC~27035 by aligning its event log with that model and prioritizing
the deviations for an auditor~\cite{palma2024compliance,palma2024benchimp}, and
support the assessor with visual analytics~\cite{palma2024eurova,palma2025impavid}.
SOC performance has also been measured empirically by injecting synthetic attacks
into an operational SOC~\cite{rosso2020saibersoc}. What remains open for the
supervisory setting is narrower: an analysis of the heterogeneous records a monitored
SOC submits (alerts, cases, investigation steps, escalations, dispositions, assets)
that yields findings citing those records, ranks organizations against each other,
keeps a human supervisor as the decision authority, and makes the decided state
verifiable afterwards---together with an honest account of how such an analysis
behaves on controlled and on independent data.

This paper studies that question with an implemented system, SAT-SA (developed for
the Smart India Hackathon problem SIH26157), and a frozen evaluation. Its emphasis is
the evidence rather than the architecture: every number is taken from a hash-verified
experiment bundle, and negative findings are reported with the same prominence as
positive ones.

\paragraph{Contributions.} We separate contribution types and state their strength.
\begin{itemize}
\item \emph{Application framing (low-confidence candidate).} Supervisor-side,
record-level analytics over SOC evidence submissions, with findings that cite
records, cross-organization peer comparison, entity ranking and a bound human
decision. Each ingredient has prior art (Section~\ref{sec:related}); the combination
is offered without a novelty claim.
\item \emph{Empirical findings.} Controlled evaluations of detection against declared
labels and simple baselines, one-worker ablation, imperfect-evidence robustness,
peer-cohort sensitivity, orchestration cost and recovery, and an external negative
result on a public IT incident log with a construct-level failure analysis.
\item \emph{Engineering.} A tenant-aware submission, analysis and review workflow with
a durable queue and a checkpointed human-review interrupt.
\item \emph{Security capability (integration).} Binding of the supervisory decision to
the analysed state with standard primitives (SHA3-256, ML-DSA-65, a hash-chained
ledger), evaluated with a mutation matrix.
\item \emph{Reproducibility practice.} Write-once experiment bundles, an evidence
freeze and a data layer that generates every reported number.
\end{itemize}
We make no claim of algorithmic novelty, superiority over existing systems, or
effectiveness in real SOCs.

The remainder of the paper reviews related work (Section~\ref{sec:related}), formalizes
the problem and the implemented analysis (Sections~\ref{sec:formulation}
and~\ref{sec:system}), describes the experimental methodology
(Section~\ref{sec:method}), reports controlled results (Section~\ref{sec:results}) and
the external evaluation (Section~\ref{sec:external}), and discusses implications,
threats to validity and reproducibility (Sections~\ref{sec:discussion}--\ref{sec:repro}).
